\section{Interpretation}

\begin{frame}{Ruling out simple alternatives}
  \small
  \begin{columns}[T,onlytextwidth]
    \begin{column}{0.52\textwidth}
      \begin{block}{Alternatives discussed in the paper}
        \begin{itemize}
          \item Incomplete or unrepresentative networks
          \item “Exceptional women” rather than a broad pattern
          \item Role monopolisation by women
          \item Over-reading a small set of centrality measures
          \item Online attention effects in male-dominated spaces
        \end{itemize}
      \end{block}
    \end{column}

    \begin{column}{0.44\textwidth}
      \begin{exampleblock}{Additional checks from the discussion}
        \begin{itemize}
          \item ISIS groups can approach 50\% women, so rarity alone is weak.
          \item In PIRA, expert reputation lists did not align with top-BC women.
          \item The paper's robustness checks extend beyond one metric or one explanation.
        \end{itemize}
      \end{exampleblock}
    \end{column}
  \end{columns}

  \note{
    \textbf{Time: 2:15}\par
    Treat this as a credibility slide.\par
    Move quickly but explicitly through the alternatives the authors themselves consider.
  }
\end{frame}

\begin{frame}{Not driven by a few exceptional women}
  \small
  \begin{columns}[T,onlytextwidth]
    \begin{column}{0.32\textwidth}
      \begin{block}{Main point}
        The pattern is broad, not driven by a tiny handful of standout women.
      \end{block}
    \end{column}

    \begin{column}{0.68\textwidth}
      \PaperFigureSlot[0.7\textheight]{S4.png}{
        Supplementary figure S4: CCDF evidence against the ``few exceptional women'' explanation.
      }
    \end{column}
  \end{columns}

  \note{
    \textbf{Time: 1:30}\par
    \textbf{MUST SAY:}\par
    The online women's betweenness distribution is broad across representative time steps rather than dominated by a few outliers.\par
    For PIRA, the highest-betweenness women did not line up with the names that subject-matter experts regarded as especially prominent.\par
    Explain the robustness point, not the CCDF mathematics in detail.\par
    Then move to the role-assignment alternative.
  }
\end{frame}

\begin{frame}{High betweenness is not tied to one operational role}
  \small
  \begin{columns}[T,onlytextwidth]
    \begin{column}{0.32\textwidth}
      \begin{block}{Main point}
        The paper finds no clear evidence that women's higher betweenness comes from occupying one special role category.
      \end{block}
    \end{column}

    \begin{column}{0.68\textwidth}
      \PaperFigureSlot[0.6\textheight]{S5.png}{
        Supplementary figure S5: normalized betweenness centrality broken down by PIRA operational role.
      }
    \end{column}
  \end{columns}

  \note{
    \textbf{Time: 1:30}\par
    \textbf{MUST SAY:}\par
    This slide rebuts the role-assignment alternative: if women systematically occupied a single high-betweenness specialty, that alone could explain the gender difference.\par
    Instead, men and women are scattered across operational roles, so the BC pattern looks more like a bottom-up network outcome than an assignment artifact.\par
    Keep the point narrow: this is a rebuttal to one alternative explanation.\par
    Then move to the paper's stated limitations.
  }
\end{frame}

\begin{frame}{Limitations and validity boundaries}
  \small
  \begin{block}{Data limitations}
    \begin{itemize}
      \item The authors cannot prove completeness; they argue the networks are representative enough to reveal real effects.
      \item External events shape network evolution and platform behaviour.
    \end{itemize}
  \end{block}

  \begin{block}{Measurement limitations}
    \begin{itemize}
      \item Online gender declarations may contain noise, though the paper argues misdeclaration incentives are limited.
      \item Higher-order network statistics need more data and are more sensitive to noise.
    \end{itemize}
  \end{block}

  {\footnotesize\color{PresTwoSecondary}\textbf{Inference discipline:}
  claims are about \alert{topological roles}, not intent or psychology.}

  \note{
    \textbf{Time: 2:15}\par
    Say clearly what the paper does not claim.\par
    Keep this slide as a boundary marker before discussing implications.
  }
\end{frame}

\begin{frame}{Implications: rethinking “importance” in networked systems}
  \small
  \begin{columns}[T,onlytextwidth]
    \begin{column}{0.55\textwidth}
      \begin{block}{Interpretive implication}
        Traditional approaches often focus on high-degree hubs; the paper argues that women's interconnectivity and \alert{collective measures like betweenness} also matter.
      \end{block}

      \begin{exampleblock}{Ethical constraint from the paper}
        The authors do \alert{not} recommend targeting women; they frame the implication as rethinking evaluation and engagement.
      \end{exampleblock}
    \end{column}

    \begin{column}{0.45\textwidth}
      \centering
      \begin{tikzpicture}[node distance=0.9cm]
        \node[draw=PresTwoMuted,rounded corners,fill=PresTwoSurface,align=center,inner sep=8pt] (hub) {Hub-centric\\\small degree};
        \node[draw=PresTwoMuted,rounded corners,fill=PresTwoSurface,align=center,inner sep=8pt,below=of hub] (broker) {Broker-centric\\\small betweenness};
        \draw[-Latex,thick] (hub) -- (broker);
      \end{tikzpicture}
    \end{column}
  \end{columns}

  \note{
    \textbf{Time: 1:30}\par
    Keep the implication conceptual, not tactical.\par
    This slide should sound like “how to evaluate importance,” not “how to exploit a network.”\par
    Then frame the model as a speculative add-on rather than core evidence.
  }
\end{frame}

\begin{frame}{A speculative mechanism: team ethic in a fission--fusion network}
  \small
  \begin{columns}[T,onlytextwidth]
    \begin{column}{0.68\textwidth}
      \PaperFigureSlot[0.65\textheight]{2E.png}{
        Show both model-vs-data panels together.
      }
    \end{column}

    \begin{column}{0.32\textwidth}
      \begin{exampleblock}{What the fit is meant to show}
        Under the paper's assumptions, a simple fission--fusion model can reproduce the broad direction of the observed PIRA trends.
      \end{exampleblock}
    \end{column}
      \end{columns}

  \note{
    \textbf{Time: 2:30}\par
    \textbf{MUST SAY:}\par
    High-BC actors are modeled as more team-oriented.\par
    That team ethic is assumed to spread within clusters containing women.\par
    The whole process runs on top of a fission--fusion network model.\par
    PIRA observations are yearly, while the proposed spreading dynamics are assumed to happen faster.\par
    Be explicit that the authors call this setup oversimplified and do not present it as causal proof.\par
    Explain fission--fusion only at a high level.\par
    Then use the supplementary schematic to make the assumptions concrete.
  }
\end{frame}

\begin{frame}{How the generative model operates}
  \small
  \begin{columns}[T,onlytextwidth]
    \begin{column}{0.70\textwidth}
      \PaperFigureSlot[0.70\textheight]{S3.png}{
        Supplementary figure S3: the generative multi-agent model used only for the paper's Fig. 2E comparison.
      }
    \end{column}

    \begin{column}{0.30\textwidth}
      \begin{block}{What the schematic specifies}
        It makes the update rules visible: coalescence versus fragmentation, kinship versus team rules, and the way team behaviour spreads across agents.
      \end{block}
    \end{column}
  \end{columns}

  \note{
    \textbf{Time: 1:30}\par
    \textbf{MUST SAY:}\par
    The schematic spells out the coalescence and fragmentation rules, the kinship-versus-team logic, and the way ``infected'' agents spread the team ethic.\par
    The supplementary text stresses that this is \alert{oversimplified} and far from a definitive mechanism; its role is to show one minimal process that can reproduce the observed trends.\par
    Use the schematic to clarify, not to over-sell the model.\par
    Then move to the takeaways.
  }
\end{frame}

\begin{frame}{Conclusions and open questions}
  \begin{block}{Three take-home messages}
    \begin{enumerate}
      \item In both the online and offline cases, women show higher \alert{betweenness centrality} despite being fewer in number.
      \item Degree and betweenness can diverge: women can avoid visible “star” status while providing brokerage.
      \item Longevity patterns are consistent with women's embeddedness having system-level consequences.
    \end{enumerate}
  \end{block}

  \begin{exampleblock}{Open questions}
    How stable is this pattern across other extreme networks? Which pressures change centrality patterns? How should contributions be evaluated: locally or collectively?
  \end{exampleblock}

  \note{
    \textbf{Time: 0:45}\par
    Deliver these as compact sentences.\par
    This is the memory slide immediately before Q\&A.
  }
\end{frame}

\begin{frame}[plain]
  \vspace{1.1cm}
  \begin{center}
    {\fontsize{34}{38}\selectfont \textbf{Questions}}
  \end{center}

  \note{
    \textbf{Time: 7:00}\par
    If questions are short, use the backup slides for definitions, null model logic, or data scales.\par
    If questions are long, stop on time and offer to share the deck afterward.
  }
\end{frame}
